\clearpage
\section{DFG v1.3 step by step}\label{app:learner-full}

This appendix gives every step of DFG v1.3 with the line numbers of the frozen file \repo{study/part10/hyb_portfolio3/dfg13.py}, which \cref{sec:learner} summarizes; the numbered equations are those of \cref{sec:round}. Parameter values are the constructor defaults, as run, and \cref{tab:params} lists them with the constants fixed in the code. \Cref{tab:dictionary,tab:params} are generated from the frozen file itself by \repo{paper_dfg/make_tables.py}, which also checks them against the theory notes \repo{study/part11/interp/theory/THEORY.md}.

% ------------------------------------------------------------------------------------------------------------------
\subsection{The dictionary}\label{app:learner-dict}

\paragraph{The linear block in the code.} In the code the linear block's inputs are ordinary degree-1 characters (line 832). They are written apart in \cref{eq:dfgclass} because they entered together, without tests of their own.

\paragraph{Sum blocks.} A block $B$ has inputs $S_B$, signs $s_i\in\{\pm1\}$, and the integer statistic
\[
t_B(x)=\sum_{i\in S_B}s_ix_i+a_B\,\tau_B(x),
\]
or its absolute value if the block is \emph{folded}. An unweighted block has $a_B=0$. A \emph{weighted} block (v1.3) has $a_B\in\{\pm2,\pm3\}$ and an inner term $\tau_B$ on inputs outside $S_B$: a literal $x_j$, a 2-parity $x_jx_l$, or an AND of two literals written as $\pm1$. The groups $G_{B,v}$ are contiguous runs of the values of $t_B$ that occur in training, each holding at least $\max(3,\lfloor0.005\,n\rfloor)$ training points. Rare values are merged into their neighbours (\code{_edges}, lines 86--97). $\mathcal V_B$ is every occupied group except the most populous one, whose indicator the constant replaces. A \emph{modulated} block multiplies its indicators by a character $\chi_{P_B}$ that is already a feature, on inputs outside $S_B$; otherwise $\chi_{P_B}=1$. A \emph{product} block $B_1\times B_2$ of two blocks with disjoint inputs uses the joint cells $G_{B_1,v_1}\times G_{B_2,v_2}$ of the pair $(t_{B_1}(x),t_{B_2}(x))$. When every value of $t_B$ is its own group, every function of $t_B$ is a combination of the constant and the block's indicators (\cref{prop:blocks}). A function of a signed sum of $k$ inputs then costs at most $k+1$ weights, where its Fourier expansion can need all $2^k$ characters on $S_B$.

\paragraph{The entries of \cref{tab:dictionary}.} ``Rows tested'': a budgeted search runs on the first $m$ rows, with $m$ set by its share of the budget (\cref{app:learner-budget}). At $d=30$ that is every row of the grid ($n\le16384$) except for degrees 5 and 6. Degree 5 sees at most 1374 rows, so it is subsampled from $n=2048$ on (687 in the final fit at $n\le512$, every row there); degree 6 sees at most 147 (73 in the final fit at $n\le512$, so it is subsampled from $n=128$ on). $N$ is the count the union bound pays for, and $\ln(4KN/\delta)$ the log term of the threshold, with $K=4$ and $\delta=0.1$. For the extensions, $b\le32$ is the number of current characters extended; for the local class, $V$ is the set of inputs of the current characters. ``Smallest $n$ that can pass'': for characters and blocks, no member can pass its test on fewer rows, for any labels (\cref{prop:floor}); this is $2c^2\ln(4KN/\delta)$ for characters and $1+2\sqrt{L_B}+2L_B$, the threshold of \cref{eq:blocktest} at $\nu=1$, for blocks. For the dense moves it is a condition of the code, not a floor of \cref{prop:floor}: the linear block's room condition $n\ge2(p+k)$, which is 62 in round 0 with all 30 inputs and lower when some degree-1 characters are already candidates (42 with 20 inputs), and $n\ge4d$ for the spectral features. For sum blocks, $N_B$ counts the supports, the sign patterns up to a global sign, the two fold variants, the $1+\mu$ choices of modulator and the $k_{\max}=24$ sizes. Every block searched in a round with $\mu$ current characters as modulators pays the factor $1+\mu$, not only the modulated ones; weighted blocks multiply $N_B$ by the factor shown, and a product pays for both factors.

\paragraph{Description length.} A fitted model's description length is the sum of its features' names; at $d=30$, par4's $\chi_{\{1,2,3,4\}}$ costs $\ln\binom{30}4=10.22$ nats to name, and $\ln(4K\binom{30}4/\delta)=15.29$ with the test's constant (\cref{sec:relates}).

% ------------------------------------------------------------------------------------------------------------------
\subsection{One round}\label{app:learner-round}

\paragraph{Residual and noise level.} The features are the columns of $\Phi\in\R^{n\times p}$, the first one constant. The weights minimize $\sum_t(y_t-(\Phi w)_t)^2+10^{-3}\,n\sum_{j\ge2}w_j^2+10^{-9}|w|^2$ (line 589). The constant is penalized only by the $10^{-9}$ term, so the residual $r=y-\Phi w$ has mean 0 up to that term. The noise level is $s=\sqrt{\max(\mathrm{mean}(r^2),10^{-12})}$.

\paragraph{The character test.} In \cref{eq:chartest}, $c=1$, $\delta=0.1$ and $K=4$ in every round (lines 602--603). $K_{\rm later}=3$ caps only the global degree after round 0. The standard error of $v_S$ is about $\sigma_m/\sqrt m$, so the test asks for $c\sqrt{2L_N}$ standard errors: a Bonferroni threshold over the class's $N$ candidates. The classes are:
\begin{itemize}
\item \emph{global}: every $\chi_S$ with $|S|=k$, for $k\le K=4$ in round 0 and $k\le K_{\rm later}=3$ after; $N=\binom dk$. Degrees 1 and 2 come from the degree-1 coefficients $c_1=X^\top r/n$ and the residual pair matrix $M=\frac1n\sum_tr_tx_tx_t^\top$. Degrees 3 and 4 come from the same matrix with the rows weighted by $x_a$ or $x_ax_b$ (lines 616--648). In round 0 of the final fit, degrees 5 and 6 ($K_{\rm hi}=6$) come from $\frac1m\sum_tr_t\chi_{T_1}(x_t)\chi_{T_2}(x_t)$ over disjoint $T_1,T_2$ of degrees $\lceil k/2\rceil$ and $\lfloor k/2\rfloor$, one product of two character tables (lines 651--663);
\item \emph{extension}: for each of the up to 32 current characters $\chi_T$ with the largest $|w_T|$, the sets $T\cup\{j\}$ ($N=bd$ for $b$ such bases) and $T\cup\{j,l\}$ ($N=b\binom d2$), from the degree-1 and pair statistics of $r\chi_T$ (lines 665--681);
\item \emph{local}: if the current characters use $2\le|V|\le12$ inputs, every $S\subseteq V$, from one Walsh--Hadamard transform of $r$ on $V$; $N=\binom{|V|}{|S|}$ (lines 683--692).
\end{itemize}
A candidate that passes in several classes keeps its largest $z_S$. The round admits the up to 64 candidates with the largest $z_S$, within the room $\min(256,\lfloor n/2\rfloor)-p$.

\paragraph{Dense moves.} The two dense moves are in lines 694--709. In \cref{eq:lintest}, $\mathrm{lin}$ is the set of inputs whose character is neither a feature nor a candidate this round, and $k=|\mathrm{lin}|$; $p$ is counted before the round's other additions. The spectral features take the eigenvectors $u$ of $\tilde M$, the symmetrized $M$ with zero diagonal and with the entries of pairs that are features or candidates set to 0, whose eigenvalues satisfy \cref{eq:spectest}, at most four of them, and only if $n\ge4d$. The first factor of \cref{eq:spectest} is the edge of the semicircle law for a $d\times d$ Wigner matrix with entries of size $s/\sqrt n$, about the size of $\tilde M$'s entries under noise. Each such $u$ adds the feature $(u\cdot x)^2-1$.

\paragraph{The sum-block test.} In \cref{eq:blocktest}, $k_{\max}=24$ (line 719). A weighted block adds $\ln\bigl(2(4d+20\binom d2)\bigr)=9.78$ nats at $d=30$ for its inner term (line 136). The count is: $a\in\{\pm2,\pm3\}$ times a literal on any input, or a 2-parity or an AND of two literals on any pair, times 2 for a separate share of $\delta$. A product pays $L_{B_1}+L_{B_2}$. A block is searched on the round's residual, but admitted only if it also passes the same threshold on $\tilde r$, the residual projected off the constant and this round's passing characters (up to 64), with $\mathrm{mean}(\tilde r^2)$ in place of $s^2$ (lines 725--738). So it must add something beyond those characters. Two kinds of product are tested once, on $\tilde r$ only, and not searched: the product of the round's two additive candidates, and the products of pairs of additive blocks admitted in earlier rounds, the latter with $\nu=(G_A-1)(G_B-1)$ for factors of $G_A$ and $G_B$ groups (lines 783--793 and 802--815). If $\mathrm{mean}(\tilde r^2)\le10^{-4}s^2$, the characters already explain the residual, and no block is searched.

% ------------------------------------------------------------------------------------------------------------------
\subsection{The block search and the block moves}\label{app:learner-search}

\paragraph{The block search.} The test does not say which $(S,s)$ to test. \code{_block_search} (lines 179--500) proposes candidates from the residual alone. It scores them by their margin, the statistic minus its threshold, on the first $\min(n,2048)$ rows. It builds its own pair matrix $M^\circ$ (line 227): the symmetrized pair matrix of the vector it searches ($M$ in a round's first additive search, $M$ of the deflated residual in the second, that of $r\chi_P$ for a modulated block), with zero diagonal and the excluded inputs' rows and columns set to 0. Unlike $\tilde M$, it keeps the entries of pairs that are features or candidates. Its degree-1 coefficients $c_{1,i}$ are those of the same vector. Its starting points are:
\begin{itemize}
\item prefixes of the inputs ordered by $|u_i|$, for the top two eigenvectors $u$ of $M^\circ$, with signs $\mathrm{sign}(u_i)$; every prefix length $k\ge2$, folded and unfolded. If $f$ is a function of $t=\sum_{i\in S}s_ix_i$, then $\E[f\,x_ix_j]=a_2s_is_j$ for distinct $i,j\in S$ and 0 for every other pair (\cref{prop:blocks}), so $M^\circ\approx a_2(ss^\top-I)$ on $S$, and its top eigenvector gives both the support and the signs;
\item prefixes of the inputs ordered by $|c_{1,i}|$, with signs $\mathrm{sign}(c_{1,i})$;
\item greedy growth from the 3 strongest pairs of $M^\circ$ (both relative signs) and the 2 strongest single inputs. Each step adds the input and sign that maximize $T_B$, until 24 inputs or 5 steps without a better margin;
\item at $n\le128$, in a round's first additive search, a beam of width 32 over $(S,s)$, started from the best pairs and scored on the first $\min(n,1024)$ rows;
\item in the additive search, v1.3's weighted candidates. These are eigenvector prefixes whose first input has weight 2 or 3; an inner term $\tau\in\{x_jx_l,\mathrm{AND}(\pm x_j,\pm x_l)\}$ on the strongest pair of $M^\circ$, with $a\in\{\pm2,\pm3\}$; and extensions $t+a\tau$ of the two best unweighted candidates. The extensions use a literal outside $S$, an input of $S$ re-weighted, or a 2-parity or AND on one of the 24 pairs outside $S$ with the largest $|M^\circ_{jl}|$.
\end{itemize}
The three best candidates are refined by single moves (add an input with either sign, remove one, flip a sign) while the margin improves, at most 72 moves. Weighted candidates are refined in the same way, with their inner term fixed. The best weighted candidate replaces the best unweighted one if its margin is larger. The winner's groups are then formed on all $n$ rows, and the winner is tested by \cref{eq:blocktest} on all $n$ rows.

\paragraph{Block moves in a round.} At $n\ge32$ (lines 715--815):
\begin{itemize}
\item up to two \emph{additive} blocks. The second is searched after the first candidate's group indicators are projected out of the residual, whether or not the first passed;
\item \emph{products}. When the round has two additive candidates, a block $B$ is searched over the joint cells with the first one, $A$, starting also from the second. When $A$ has at least two inputs of each sign, a \emph{split} is also tried: $B$ is searched over the joint cells with the block of $A$'s positive inputs, starting from its negative inputs. The better product is admitted only if its margin on $\tilde r$ exceeds $A$'s. Otherwise the product of the round's two additive candidates is tested, if their inputs are disjoint;
\item one \emph{modulated} block for each of the up to 4 current characters $\chi_P$ with the largest $|w_P|$, searched on $r\chi_P$ with $P$'s inputs excluded;
\item the products of pairs of additive blocks, admitted in earlier rounds, that have disjoint inputs.
\end{itemize}

% ------------------------------------------------------------------------------------------------------------------
\subsection{Refit, stopping and readout}\label{app:learner-stop}

\paragraph{Refit and stopping.} The round adds the admitted characters, the linear block's inputs, each admitted block's indicators (if they fit in the room) and the spectral features, and refits the weights. DFG stops (lines 596--598, 824):
\begin{itemize}
\item when a round admits nothing (in the default variant; \cref{sec:smalln} for the forced one);
\item when no room is left, $p\ge\min(256,\lfloor n/2\rfloor)$;
\item from round 1 on, when the fit has every training label's sign right;
\item at the start of a round, if more than 0.75 of the call's budget is spent;
\item after 8 rounds.
\end{itemize}
The code also stops if $s<10^{-6}$. That cannot happen, since $s\ge10^{-6}$ by its definition (\repo{study/part11/interp/calib/VERIFY_CALIB.md}).

\paragraph{Readout.} After the last round, if $n\ge8p$, the least-squares weights are refined by 15 Newton steps on the mean logistic loss plus $\frac12\cdot10^{-4}|w|^2$, started from $4w$. The result is rescaled so that its output on the training rows has the norm of the least-squares output (lines 849--855). Part 6 added this step after the least-squares readout plateaued at 5--7\% test error on a dense halfspace (thr30) where the networks reached 1\% \cite{discretelearners2026}. The prediction adds the kernel background $g\cdot x$ (line 856).

% ------------------------------------------------------------------------------------------------------------------
\subsection{Small samples}\label{app:learner-smalln}

\paragraph{The forced variant.} With few points the union bound passes nothing. In round 0, where only the global classes run, no character can pass its test at $d=30$ below $n=17$, the degree-1 floor $2c^2\ln(4Kd/\delta)=16.95$, and no degree-4 character below $n=31$ (\cref{tab:dictionary}). Later rounds have smaller classes and lower floors. In a round where nothing passes, the forced variant adds a guess anyway, at most twice per fit (lines 817--822). If $n\ge48$ and the round's first additive block candidate has $T_B$/threshold above the square of the best character's $z$, it forces that block; both ratios are on a $\chi^2$ scale. Otherwise it adds the character with the largest $z$ among those offered. These are every degree-1 character, the largest pair entry, the largest entry of each row block of degrees 3 and 4, and, from round 1 on, the local subsets above the smallest local threshold (\repo{study/part11/interp/theory/THEORY.md}, \S0). A forced guess passes no test. Forcing a block is on by default (\code{sum_force=True}); the docstring's ``off by default'' (line 73) is out of date.

\paragraph{The paired cross-validation.} At $n\le\code{cv_n}=512$, \code{DFG.fit} chooses between the default and the forced variant by cross-validation (lines 548--571). There are $F=10$ folds at $n\le64$ and 5 above, assigned by a random permutation of the rows. Each variant is fitted on each fold's training rows, without the degree-5/6 search, and predicts the held-out fold. Let $e_{10}$ be the number of points where the default errs and the forced variant does not, and $e_{01}$ the reverse. The forced variant is chosen if $e_{10}-e_{01}>2\sqrt{e_{10}+e_{01}}$: McNemar's test for paired errors \cite{mcnemar1947note}, one-sided at two standard errors. The chosen variant is then fitted on all $n$ rows, and the fold fits' FLOPs count toward the fit. DFG's own choices on the atlas's cells are counted in \cref{app:interp-fits}. No level is proved for this choice, because the folds' held-out errors are dependent (\cref{sec:guarantees}).

\paragraph{The kernel background and the room.} The output adds $g\cdot x$, with $g=X^\top r/n$ on the final residual. With no feature but the constant, $r=y-w_0$, and the predictor is $\mathrm{sign}(w_0+g\cdot x)$: the centred Hebbian classifier, one step on all $d$ inputs, the counterpart of the networks' near-kernel component. Once the features fit the labels, $r\approx0$ and $g\approx0$. The background has no fitted parameter and costs $2nd$ FLOPs. Least squares, likewise, is kept away from its interpolation peak: characters and block indicators are admitted only within $p\le\min(256,\lfloor n/2\rfloor)$, and the linear block only if $n\ge2(p+k)$, with $p$ counted before the round's other additions. The linear block's inputs and the spectral features are appended without a room check (lines 826--844), so $p$ can exceed $n/2$ \status{Proved, by the code}; in none of the 407 atlas cells does it exceed $\min(256,\lfloor n/2\rfloor)$ (\repo{study/part11/interp/tables/atlas_cells.csv}) \status{Numerical}.

% ------------------------------------------------------------------------------------------------------------------
\subsection{The budget}\label{app:learner-budget}

\paragraph{The budget.} The unit is one training run of a toy network at the same $d$, $\mathrm{GD}(d)=394\,PT$ FLOPs for $P=128d+16768$ parameters and $T$ steps: $1.62\times10^{11}$ at $d=30$ (\cref{sec:metrics}). A fit call gets $\mathrm{GD}\cdot\text{scale}$, with scale 1 at $n>512$. At $n\le512$ the final fit gets $\frac12$ and each of the $2F$ fold fits $1/(4F)$, so the shares add up to one run. Within a call:
\begin{itemize}
\item the exhaustive searches run on the first $m$ rows, with $m$ set so that each costs at most a fixed share of the call's budget. The shares are 15\% for degree 4 (round 0 only), 3\% per round for degree 3, 3\% per round for the two-input extensions, and 3\% each for degrees 5 and 6 (round 0 of the final fit only). The search uses $m=\min(n,\lfloor\text{share}\cdot\mathrm{GD}\cdot\text{scale}/c_{\rm row}\rfloor)$ rows, with $c_{\rm row}$ the code's per-row cost estimate (\cref{prop:compute}), and its threshold uses $m$ (\cref{eq:chartest}). A search whose $m$ falls below 32 rows (64 for degrees 5 and 6) is skipped. At $d=30$ only degrees 5 and 6 are subsampled: degree 5 to 1374 rows (687 in the final fit at $n\le512$, every row there), degree 6 to 147 (73 in the final fit at $n\le512$). At $d=60$ the degree-4 search sees at most 2212 rows, and at $d=120$ at most 364, with that $d$'s 40000-step run as the unit;
\item a round starts only while the call has spent less than 0.75 of its budget. The other parts are not budgeted: degrees 1 and 2, the one-input extensions, the local class, the eigendecomposition, the sum blocks, least squares, the Newton steps and the background. Only this cap bounds them, so a call can exceed 0.75 by the cost of its last round.
\end{itemize}

\begin{table}[t]
\centering{\scriptsize\setlength{\tabcolsep}{3pt}
\begin{minipage}[t]{0.505\linewidth}
\begin{tabularx}{\linewidth}[t]{@{}l@{\hspace{4pt}}l@{\hspace{5pt}}>{\raggedright\arraybackslash}X@{}}
\toprule
\multicolumn{3}{@{}l}{\emph{Tests}}\\
\code{K} & 4 & global degree in round 0\\
\code{K_later} & 3 & global degree, rounds $\ge1$\\
\code{K_hi} & 6 & top degree, round 0\\
\code{c} & 1 & threshold multiplier\\
\code{delta} & 0.1 & level $\delta$ of every test\\
\code{noise} & \code{residual} & tests scaled by $s=\mathrm{rms}(r)$\\
\multicolumn{3}{@{}l}{\emph{Rounds and classes}}\\
\code{rounds} & 8 & rounds at most\\
\code{max_new} & 64 & characters admitted per round\\
\code{max_feats} & 256 & features (also $\le n/2$)\\
\code{max_ext_bases} & 32 & characters extended\\
\code{local_max} & 12 & largest $|V|$, local class\\
\code{blocks} & True & linear and spectral moves\\
\multicolumn{3}{@{}l}{\emph{Small $n$ and budget}}\\
\code{small_n} & True & forced variant and its test\\
\code{cv_n} & 512 & the test runs at $n\le$ this\\
\code{background} & True & kernel background $g\cdot x$\\
\code{budget} & (0.15, 0.03, 0.03) & shares: degree 4; 3; ext.\ by 2\\
\code{budget_hi} & 0.03 & share: degrees 5, 6 each\\
\multicolumn{3}{@{}l}{\emph{Readout}}\\
\code{readout} & \code{ls} & run as \code{logit}\\
\code{ridge} & 0.001 & ridge of least squares ($\times n$)\\
\code{logit_lam} & 0.0001 & L2 penalty, logistic steps\\
\bottomrule
\end{tabularx}
\end{minipage}\hfill
\begin{minipage}[t]{0.47\linewidth}
\begin{tabularx}{\linewidth}[t]{@{}l@{\hspace{4pt}}l@{\hspace{5pt}}>{\raggedright\arraybackslash}X@{}}
\toprule
\multicolumn{3}{@{}l}{\emph{Sum blocks}}\\
\code{sums} & True & sum blocks\\
\code{sum_nmin} & 32 & smallest $n$ for the block search\\
\code{sum_kmax} & 24 & largest block; the $k_{\max}$ of $N_B$\\
\code{sum_mod} & 4 & modulators $\mu$ at most\\
\code{sum_add} & 2 & additive blocks per round\\
\code{sum_gcap} & 2048 & rows for the search's scores\\
\code{sum_mc} & 0.005 & smallest group over $n$ (at least 3)\\
\code{sum_prod} & True & product and split blocks\\
\code{sum_seeds} & (2, 3, 2) & eigenvector, pair, single seeds\\
\code{sum_inner} & True & weighted blocks $t+a\tau$\\
\code{sum_ntau} & 1 & pairs tried as inner terms\\
\code{sum_beam} & 32 & beam width\\
\code{sum_beam_nmax} & 128 & beam only at $n\le$ this\\
\code{sum_force} & True & forced variant may force a block\\
\code{sum_force_nmin} & 48 & smallest $n$ for a forced block\\
\code{sum_force_r2} & 0 & smallest $T_B/n$ of a forced block\\
\bottomrule
\end{tabularx}
\par\medskip
\begin{tabularx}{\linewidth}[t]{@{}l@{\hspace{5pt}}>{\raggedright\arraybackslash}Xr@{}}
\toprule
\multicolumn{2}{@{}l}{\emph{Fixed in the code}} & line\\
15 & logistic Newton steps & 851\\
$n\ge8p$ & when they run & 849\\
$4w$ & where they start & 850\\
0.75 & budget share that ends the rounds & 598\\
10 / 5 & folds at $n\le64$ / above & 553\\
$2\sqrt{e_{10}+e_{01}}$ & margin of the paired test & 567\\
$\frac1{4F}$, $\frac12$ & budget scale: fold fit, final fit & 564, 569\\
32 / 64 & fewest rows: degrees 3, 4 / 5, 6 & 630, 655\\
4; $4d$ & spectral features; smallest $n$ & 707\\
2 & forced guesses per fit & 817\\
$10^{-4}$ & least $\mathrm{mean}(\tilde r^2)/s^2$ for blocks & 734\\
3 & fewest points per group & 721\\
$\pm2,\pm3$ & weights $a$ of inner terms & 273\\
5 & growth steps without gain & 179\\
24 & pairs for inner-term extensions & 180\\
1024 & rows of the beam & 180\\
$3k_{\max}$ & refinement moves at most & 376\\
\bottomrule
\end{tabularx}
\end{minipage}
}
\caption{DFG v1.3's constructor defaults, as run (\code{DFG(readout="logit")}, every other argument at its default), and constants fixed in its code, with their lines \status{Proved, by the code}. Source: \repo{study/part10/hyb_portfolio3/dfg13.py}, read by \repo{paper_dfg/make_tables.py} (\texttt{tab\_params}). The defaults are read from the constructor's signature, and each fixed constant is checked to occur at the line given.}
\label{tab:params}
\end{table}
